\section{Discussion}\label{sec:discussion}

Our findings suggest that executable architecture rules function more like governance mechanisms than static design documents. They are adopted selectively, maintained irregularly, and can become permissive through exemptions and ignore lists. At the same time, when enforced in CI and maintained with discipline, they can improve architectural coherence and constrain drift. This duality is the central implication of the study: AaC is both a technical practice and an organizational discipline.

An exploratory but important question concerns the role of automated agents in architectural governance. As AI-assisted development becomes more common, architecture rules may serve as guardrails against agent-generated drift: large language models and coding agents can quickly generate modules, imports, or cross-cutting references that violate intended boundaries. Executable checks, especially when enforced in CI, can act as protective constraints that detect violations early and reduce the need for downstream architecture repair. This raises a new research question: to what extent can architecture-as-code serve as a control mechanism for AI-induced architecture regressions? We treat this as a motivating exploratory direction for future work, not a final claim.

The implications are actionable across multiple stakeholder groups. For researchers, the benchmark and longitudinal evidence provide a foundation for studying rule drift, governance quality, and architecture maintenance under real-world constraints. For developers, the results highlight the need to treat exemptions as a governance signal rather than a harmless convenience. For architecture framework designers, the evidence suggests that future AaC systems should incorporate explicit awareness of rule rot, enforcement quality, and causal monitoring in addition to syntactic validation. This aligns with broader calls for continuous architectural governance and conformance-aware development infrastructures~\cite{contaac,archguard}.

